\section{从 John 到 Sandy：工作流中的每一步}
\subsection{一次普通的故障单：人工流程的长链}
在 00:02:11--00:11:21 的故事里，John 因为早晨咖啡洒在笔记本上，向 IT 提单请求补发设备；该工单分配给 Sandy，Sandy 逐条查阅知识库、历史事件，再调整权限并写清记录。即便最新的生成式 AI 已经部署，Sandy 依然执行大量手工步骤，包括核对访问控制、召唤相关系统、写解决方案摘要。

\begin{warningbox}{自动化的“多米诺骨牌”缺口}
即便有 AI 摘要，企业工作流中仍有数百种场景需要人工判断、权限核查、合规判断与多系统协调，任何自动化都需要在此过程中保持可追溯与人工干预的余地。
\end{warningbox}

\subsection{低频但关键的“沙粒任务”}
讲者指出，大多数“低价值、低频”任务都由于过于独特而传统脚本难以覆盖，这些“沙粒任务”每天在全球企业中重复百万次，但每一次都需要快速响应。

\begin{importantbox}{Agent 的覆盖面不该只看“高频”
}
Agent 应该抓住“可解释、可追溯、可细粒度调整”的机会：即使单次任务耗时短，但在 aggregate 级别掀起的效率提升和员工满意度才是关键。
\end{importantbox}

\subsection{本章小结}
John/Sandy 的例子展现了 Agent 在多系统、多角色、多步骤场景中的潜在价值与风险点，为后文的架构和审计需求埋下伏笔。

